The Python files in \texttt{source code/} are educational demonstrations. The maintained command-line examples are \texttt{tfhe\_example.py}, \texttt{bfv\_example.py}, \texttt{ckks\_example.py}, and \texttt{bgv\_example.py}; the README lists their supported operations.

The upstream source is available at \url{https://github.com/fhetextbook/fhetextbook.github.io}, and the reviewed bilingual fork at \url{https://github.com/wang2000star/fhe-textbook-cn}.

The implemented operations, simulations, and missing functions must be distinguished. TFHE's \texttt{tfhe\_mock\_pbs} reads a secret key, decrypts, applies a lookup table, and re-encrypts. Its ``bootstrapping key'' dictionary contains that secret key. It is a client-side functional simulation, not a public evaluation key or a secure server-side implementation of programmable/circuit bootstrapping. The default ciphertext multiplication example instead demonstrates a tensor-product and key-switching construction and does not benchmark standard TFHE bootstrapping. BFV, CKKS, and BGV bootstrapping functions are unimplemented and explicitly raise \texttt{NotImplementedError}.

Small toy dimensions, bounded sample errors, and Python/NumPy pseudorandom sampling are used to illustrate algebra. They are not a validated cryptographic parameter set or a production random-number generator. Successful example runs verify only those chosen inputs and do not establish a security level, arbitrary-depth correctness, or the claimed performance of an optimized FHE library. The older standalone encoding and RNS scripts are algebra experiments; see the academic review for their coverage and limitations.
